%
% reversi_290_nuovi_pattern.tex
%
% This file is part of the reversi program
% http://github.com/rcrr/reversi-v2
%
% Author Roberto Corradini mailto:rob_corradini@yahoo.it
% Copyright 2026 Roberto Corradini. All rights reserved.
%
% License
%
% This program is free software; you can redistribute it and/or modify it
% under the terms of the GNU General Public License as published by the
% Free Software Foundation; either version 3, or (at your option) any
% later version.
%
% This program is distributed in the hope that it will be useful,
% but WITHOUT ANY WARRANTY; without even the implied warranty of
% MERCHANTABILITY or FITNESS FOR A PARTICULAR PURPOSE. See the
% GNU General Public License for more details.
%
% You should have received a copy of the GNU General Public License
% along with this program; if not, write to the Free Software
% Foundation, Inc., 59 Temple Place - Suite 330, Boston, MA  02111-1307, USA
% or visit the site <http://www.gnu.org/licenses/>.
%

%-----------------------------------------------------------------------------
% Capitolo: Nuovi pattern.
%
% Tre blocchi di esperimenti a ec20 su pattern nuovi o estesi.
%
% 1) Il pattern CASTLE (12 bit) al posto di XEDGE (10 bit), in tre run. A
%    parita di condizioni (batch 3/6, min_grad [1e-05, 3]) XEDGE resta
%    avanti di 0.09 in validazione; con i batch 3,9,12 e il criterio
%    [1e-06, 3] CASTLE ottiene 12.51, 0.40 sotto il riferimento XEDGE a
%    batch singolo. Nel feature set completo del modello 0108, CASTLE al
%    posto di XEDGE porta la RMSE di validazione da 7.01 a 6.93 (il primo
%    modello della serie sotto 7.00) e i pesi attivi da 81.269 a 219.546.
%
% 2) Il pattern 2X6COR (12 bit) al posto di 2X5COR (10 bit), in tre run. A
%    parita di condizioni 2X5COR resta avanti di 0.07 in validazione; con i
%    batch 3,9,12 2X6COR ottiene 9.08, 0.59 sotto il riferimento 2X5COR a
%    batch singolo. Nel feature set completo del modello 0108, 2X6COR al
%    posto di 2X5COR porta la RMSE di validazione da 7.01 a 6.85 (il primo
%    modello della serie sotto 6.93).
%
% 3) Cinque pattern nuovi (3X4COR, KATE, KATE2, KATE3, KATEK) in tre run
%    ciascuno, al posto di CORNER. Tutti e cinque battono CORNER da soli a
%    parita di condizioni e nel feature set completo del modello 0108;
%    3X4COR e' il migliore nelle tre condizioni (6.78, il nuovo record
%    della serie, 0.07 sotto il 2X6COR).
%
% 4) Il trio: il feature set del 0108 con XEDGE sostituito da CASTLE, 2X5COR
%    da 2X6COR e CORNER da KATEK (modello 0130): 6.68 in validazione, 0.33
%    sotto il 0108, 0.07 oltre la somma degli effetti misurati singolarmente.
%    Sul trio, i cinque corner a confronto (0130-0134): l'ordinamento e'
%    identico a quello del contesto 0108; 3X4COR e' il nuovo campione della
%    serie (0134, 6.57), con 670.519 pesi attivi e il gap piu' alto (0.27).
%
% 5) Sei pattern piu' semplici (RCT2X4, FOURC, CORE, CORED, COREA, BARBEL):
%    i cinque centrali non hanno segnale a ec20 (da soli RMSE 22-26, 7-349
%    pesi attivi; nel modello completo sono inerti, 0.00). BARBEL ha
%    segnale (15.43/15.24) e la sua sostituzione di DIAG8 nel 0134 e'
%    l'intero miglioramento del modello a 15 pattern (0147, 6.49): il
%    modello 0153 (0134 con BARBEL al posto di DIAG8, 6.50) diventa il
%    nuovo riferimento della serie.
%-----------------------------------------------------------------------------

\chapter{Nuovi pattern}

\section{Da XEDGE a CASTLE}

Il pattern CASTLE (mask \texttt{0x000000000000C3FF}, 12 bit)
estende il pattern XEDGE
(\texttt{0x00000000000042FF}, 10 bit). Tre run a ec20 lo
mettono alla prova: il modello 0109 (intercept + CASTLE, batch
3/6, \texttt{min\_grad} $[1\times10^{-5}, 3]$), eseguito alle
medesime condizioni del modello 0053 (intercept + XEDGE); il
modello 0110 (medesimo feature set, batch 3, 9, 12/6,
\texttt{min\_grad} $[1\times10^{-6}, 3]$, \texttt{max\_iters}
2000); il modello 0111 (il feature set del modello 0108 con
CASTLE al posto di XEDGE: dieci pattern e la mobilità OptMob,
configurazione di riferimento, \texttt{min\_grad}
$[1\times10^{-6}, 5]$). La tabella~\ref{tab:castle} riassume le
tre run CASTLE e i due modelli XEDGE di confronto.

\begin{table}[ht]
\centering
\small
\begin{tabular}{llllrrrr}
\toprule
modello & feature set & batch train & \texttt{min\_grad} & RMSE train & RMSE val & $W$ & iterazioni \\
\midrule
0053 & intercept + XEDGE & 3 & $[10^{-5}, 3]$ & 12.78 & 12.91 & 24\,047 & 171 \\
0109 & intercept + CASTLE & 3 & $[10^{-5}, 3]$ & 11.93 & 13.00 & 141\,032 & 284 \\
0110 & intercept + CASTLE & 3, 9, 12 & $[10^{-6}, 3]$ & 12.32 & 12.51 & 164\,239 & 819 \\
0108 & 10 pattern + XEDGE + OptMob & 3, 9, 12 & $[10^{-6}, 5]$ & 6.95 & 7.01 & 81\,269 & 687 \\
0111 & 10 pattern + CASTLE + OptMob & 3, 9, 12 & $[10^{-6}, 5]$ & 6.82 & 6.93 & 219\,546 & 909 \\
\bottomrule
\end{tabular}
\caption{Il pattern CASTLE in tre run a ec20, accanto ai modelli
XEDGE di confronto (0053 e 0108).}
\label{tab:castle}
\end{table}

A parità di condizioni (0053 contro 0109) la RMSE di validazione
di XEDGE è 0.09 più bassa di quella di CASTLE (12.91 contro
13.00); CASTLE porta 141\,032 pesi attivi contro i 24\,047 di
XEDGE e una RMSE di training più bassa (11.93 contro 12.78). Con
i batch 3, 9 e 12 di training e il criterio
$[1\times10^{-6}, 3]$ (0110) la RMSE di validazione di CASTLE è
12.51, 0.40 sotto il riferimento XEDGE a batch singolo. Nel
feature set completo (0111 contro 0108) la sostituzione di
XEDGE con CASTLE abbassa la RMSE di validazione da 7.01 a 6.93
e quella di training da 6.95 a 6.82; i pesi attivi passano da
81\,269 a 219\,546 (+138\,277) e il gap training--validazione da
0.06 a 0.11. Entrambe le run convergono (687 e 909
iterazioni). Il modello 0111 è il primo della serie con una RMSE
di validazione inferiore a 7.00.

\section{Da 2X5COR a 2X6COR}

Il pattern 2X6COR (mask \texttt{0x0000000000003F3F}, 12 bit)
estende il pattern 2X5COR
(\texttt{0x0000000000001F1F}, 10 bit). Tre run a ec20 lo
mettono alla prova: il modello 0112 (intercept + 2X6COR, batch
3/6, \texttt{min\_grad} $[1\times10^{-5}, 3]$), eseguito alle
medesime condizioni del modello 0054 (intercept + 2X5COR); il
modello 0113 (medesimo feature set, batch 3, 9, 12/6,
\texttt{min\_grad} $[1\times10^{-6}, 3]$, \texttt{max\_iters}
2000); il modello 0114 (il feature set del modello 0108 con
2X6COR al posto di 2X5COR: dieci pattern e la mobilità
OptMob, configurazione di riferimento, \texttt{min\_grad}
$[1\times10^{-6}, 5]$). La tabella~\ref{tab:2x6cor} riassume le
tre run 2X6COR e i due modelli 2X5COR di confronto.

\begin{table}[ht]
\centering
\small
\begin{tabular}{llllrrrr}
\toprule
modello & feature set & batch train & \texttt{min\_grad} & RMSE train & RMSE val & $W$ & iterazioni \\
\midrule
0054 & intercept + 2X5COR & 3 & $[10^{-5}, 3]$ & 9.45 & 9.67 & 37\,848 & 139 \\
0112 & intercept + 2X6COR & 3 & $[10^{-5}, 3]$ & 8.28 & 9.74 & 259\,768 & 273 \\
0113 & intercept + 2X6COR & 3, 9, 12 & $[10^{-6}, 3]$ & 8.80 & 9.08 & 306\,009 & 867 \\
0108 & 10 pattern + 2X5COR + OptMob & 3, 9, 12 & $[10^{-6}, 5]$ & 6.95 & 7.01 & 81\,269 & 687 \\
0114 & 10 pattern + 2X6COR + OptMob & 3, 9, 12 & $[10^{-6}, 5]$ & 6.70 & 6.85 & 348\,122 & 914 \\
\bottomrule
\end{tabular}
\caption{Il pattern 2X6COR in tre run a ec20, accanto ai modelli
2X5COR di confronto (0054 e 0108).}
\label{tab:2x6cor}
\end{table}

A parità di condizioni (0054 contro 0112) la RMSE di
validazione di 2X5COR è 0.07 più bassa di quella di 2X6COR
(9.67 contro 9.74); 2X6COR porta 259\,768 pesi attivi contro i
37\,848 di 2X5COR e una RMSE di training molto più bassa
(8.28 contro 9.45). Con i batch 3, 9 e 12 di training e il
criterio $[1\times10^{-6}, 3]$ (0113) la RMSE di validazione
di 2X6COR è 9.08, 0.59 sotto il riferimento 2X5COR a batch
singolo. Nel feature set completo (0114 contro 0108) la
sostituzione di 2X5COR con 2X6COR abbassa la RMSE di
validazione da 7.01 a 6.85 e quella di training da 6.95 a
6.70; i pesi attivi passano da 81\,269 a 348\,122
(+266\,853) e il gap training--validazione da 0.06 a 0.15.
Entrambe le run convergono (687 e 914 iterazioni). Il modello
0114 è il primo della serie con una RMSE di validazione
inferiore a 6.93.

\section{Cinque alternative a CORNER}

Cinque nuovi pattern entrano nel catalogo (appendice~A):
3X4COR (12 bit), KATE (10), KATE2 (11), KATE3 (12) e KATEK
(12). Ognuno è messo alla prova con il protocollo a tre run
già usato per CASTLE e 2X6COR: da solo con batch 3/6 e
\texttt{min\_grad} $[1\times10^{-5}, 3]$, contro il modello
0052 (intercept + CORNER, medesime condizioni); da solo con i
batch 3, 9, 12 e il criterio $[1\times10^{-6}, 3]$; e nel
feature set completo del modello 0108 al posto di CORNER
(configurazione di riferimento, \texttt{min\_grad}
$[1\times10^{-6}, 5]$). I quindici modelli sono i
0115--0129, tre per pattern. La tabella~\ref{tab:corner5} li
riassume accanto ai due riferimenti 0052 e 0108.

\begin{table}[ht]
\centering
\small
\begin{tabular}{llllrrrr}
\toprule
modello & feature set & batch train & \texttt{min\_grad} & RMSE train & RMSE val & $W$ & iterazioni \\
\midrule
0052 & intercept + CORNER & 3 & $[10^{-5}, 3]$ & 12.21 & 12.24 & 5\,157 & 109 \\
0115 & intercept + 3X4COR & 3 & $[10^{-5}, 3]$ & 9.02 & 9.94 & 163\,043 & 222 \\
0118 & intercept + KATE & 3 & $[10^{-5}, 3]$ & 12.05 & 12.11 & 10\,192 & 153 \\
0121 & intercept + KATE2 & 3 & $[10^{-5}, 3]$ & 11.99 & 12.11 & 20\,047 & 165 \\
0124 & intercept + KATE3 & 3 & $[10^{-5}, 3]$ & 11.79 & 12.21 & 56\,657 & 148 \\
0127 & intercept + KATEK & 3 & $[10^{-5}, 3]$ & 10.69 & 11.16 & 68\,725 & 235 \\
0116 & intercept + 3X4COR & 3, 9, 12 & $[10^{-6}, 3]$ & 9.36 & 9.54 & 189\,347 & 776 \\
0119 & intercept + KATE & 3, 9, 12 & $[10^{-6}, 3]$ & 12.08 & 12.07 & 10\,418 & 352 \\
0122 & intercept + KATE2 & 3, 9, 12 & $[10^{-6}, 3]$ & 12.04 & 12.04 & 20\,725 & 455 \\
0125 & intercept + KATE3 & 3, 9, 12 & $[10^{-6}, 3]$ & 11.95 & 12.01 & 61\,015 & 454 \\
0128 & intercept + KATEK & 3, 9, 12 & $[10^{-6}, 3]$ & 10.85 & 10.93 & 75\,155 & 686 \\
0108 & 10 pattern + CORNER + OptMob & 3, 9, 12 & $[10^{-6}, 5]$ & 6.95 & 7.01 & 81\,269 & 687 \\
0117 & 10 pattern + 3X4COR + OptMob & 3, 9, 12 & $[10^{-6}, 5]$ & 6.64 & 6.78 & 265\,389 & 856 \\
0120 & 10 pattern + KATE + OptMob & 3, 9, 12 & $[10^{-6}, 5]$ & 6.93 & 6.99 & 86\,460 & 684 \\
0123 & 10 pattern + KATE2 + OptMob & 3, 9, 12 & $[10^{-6}, 5]$ & 6.90 & 6.97 & 96\,767 & 750 \\
0126 & 10 pattern + KATE3 + OptMob & 3, 9, 12 & $[10^{-6}, 5]$ & 6.86 & 6.95 & 137\,057 & 774 \\
0129 & 10 pattern + KATEK + OptMob & 3, 9, 12 & $[10^{-6}, 5]$ & 6.79 & 6.88 & 151\,197 & 842 \\
\bottomrule
\end{tabular}
\caption{I cinque pattern alternativi a CORNER in tre run a ec20
ciascuno, accanto ai riferimenti 0052 (CORNER da solo) e 0108
(feature set completo con CORNER).}
\label{tab:corner5}
\end{table}

Da soli a parità di condizioni (batch 3/6) tutti e cinque i
pattern battono CORNER (12.24): 3X4COR di 2.30 (9.94), KATEK
di 1.08 (11.16), KATE, KATE2 e KATE3 di 0.13, 0.13 e 0.03
(12.11, 12.11, 12.21). Con i batch 3, 9 e 12 la RMSE di
validazione di 3X4COR scende a 9.54, sotto il riferimento
2X5COR a batch singolo (9.67 del modello 0054); KATEK è
secondo a 10.93.

Nel feature set completo tutti e cinque battono il modello
0108 (7.01): 3X4COR porta la RMSE di validazione a 6.78,
0.23 sotto e il nuovo record della serie, 0.07 sotto il 2X6COR
del modello 0114 (6.85); seguono KATEK a 6.88, KATE3 a 6.95,
KATE2 a 6.97 e KATE a 6.99. I quindici run convergono tutti.

3X4COR è il pattern migliore della campagna in tutte e tre le
condizioni. KATE ha il profilo di generalizzazione più pulito
(gap training--validazione tra 0.00 e 0.06 nelle tre run) e,
nel modello completo, i pesi attivi più bassi (86\,460). Il
gap del modello 0117 è 0.14, nella stessa fascia del 0114
(0.15).

\section{Il trio: CASTLE, 2X6COR e il corner}

Il modello 0130 prende il feature set del modello 0108 e
applica le tre sostituzioni già misurate nei paragrafi
precedenti: XEDGE con CASTLE, 2X5COR con 2X6COR e CORNER con
KATEK — dieci pattern e la mobilità OptMob, configurazione di
riferimento, \texttt{min\_grad} $[1\times10^{-6}, 5]$. La run
converge in 972 iterazioni: RMSE di training 6.45, di
validazione 6.68, 556\,327 pesi attivi. Le tre sostituzioni
valgono da sole 0.02, 0.08 e 0.16 nel contesto 0108
(tabella~\ref{tab:corner5}); insieme valgono 0.33, cioè 0.07
oltre la somma.

Sul trio si confrontano i cinque corner: il modello 0130
(KATEK) e i cloni 0131--0134, identici in ogni altro campo,
con KATE, KATE2, KATE3 e 3X4COR al posto di KATEK. La
tabella~\ref{tab:trio} li riassume.

\begin{table}[ht]
\centering
\small
\begin{tabular}{llrrrrr}
\toprule
modello & corner & RMSE train & RMSE val & $W$ & iterazioni & gap \\
\midrule
0134 & 3X4COR & 6.30 & 6.57 & 670\,519 & 944 & 0.27 \\
0130 & KATEK & 6.45 & 6.68 & 556\,327 & 972 & 0.23 \\
0133 & KATE3 & 6.52 & 6.75 & 542\,187 & 945 & 0.23 \\
0132 & KATE2 & 6.56 & 6.77 & 501\,897 & 1002 & 0.21 \\
0131 & KATE & 6.58 & 6.79 & 491\,590 & 983 & 0.21 \\
\bottomrule
\end{tabular}
\caption{Il trio (2X6COR + CASTLE + corner) a ec20 con i cinque
corner; il resto del feature set è quello del modello 0108.}
\label{tab:trio}
\end{table}

L'ordinamento dei cinque corner (3X4COR, KATEK, KATE3, KATE2,
KATE) è identico a quello misurato nel contesto 0108
(tabella~\ref{tab:corner5}): i passi sono 0.11, 0.07, 0.02 e
0.02, e KATE, KATE2 e KATE3 formano un plateau. Il modello
0134 (3X4COR) è il nuovo campione della serie: 6.57 in
validazione, 0.44 sotto il modello 0108 e 0.11 sotto il 0130.
I pesi attivi crescono col corner, da 491\,590 (KATE) a
670\,519 (3X4COR), contro gli 81\,269 del 0108; il gap
training--validazione va da 0.21 a 0.27, massimo per il 0134.
Le cinque run convergono (944--1002 iterazioni).

\section{I pattern semplici}

Sei pattern più semplici sono messi alla prova con il protocollo
a due run: da soli (intercept + pattern), batch 3/6 e \texttt{min\_grad}
$[1\times10^{-5}, 3]$, contro il modello 0052 (CORNER, 12.24); da
soli con i batch 3, 9, 12 e il criterio $[1\times10^{-6}, 3]$. La
tabella~\ref{tab:semplici} li riassume.

\begin{table}[ht]
\centering
\small
\begin{tabular}{lllrrrr}
\toprule
modello & pattern & batch train & RMSE train & RMSE val & $W$ & iterazioni \\
\midrule
0052 & CORNER & 3 & 12.21 & 12.24 & 5\,157 & 109 \\
0135 & RCT2X4 & 3 & 25.74 & 25.76 & 349 & 41 \\
0136 & RCT2X4 & 3, 9, 12 & 25.76 & 25.76 & 351 & 117 \\
0137 & FOURC & 3 & 22.22 & 22.23 & 22 & 33 \\
0138 & FOURC & 3, 9, 12 & 22.23 & 22.23 & 22 & 43 \\
0139 & CORE & 3 & 26.40 & 26.39 & 7 & 11 \\
0140 & CORE & 3, 9, 12 & 26.42 & 26.39 & 7 & 12 \\
0141 & CORED & 3 & 26.21 & 26.21 & 232 & 38 \\
0142 & CORED & 3, 9, 12 & 26.22 & 26.20 & 232 & 67 \\
0143 & COREA & 3 & 25.77 & 25.80 & 331 & 40 \\
0144 & COREA & 3, 9, 12 & 25.79 & 25.80 & 333 & 61 \\
0145 & BARBEL & 3 & 15.10 & 15.43 & 37\,317 & 189 \\
0146 & BARBEL & 3, 9, 12 & 15.23 & 15.24 & 42\,818 & 768 \\
\bottomrule
\end{tabular}
\caption{I sei pattern semplici da soli a ec20, accanto al
riferimento 0052 (CORNER).}
\label{tab:semplici}
\end{table}

I cinque pattern centrali (RCT2X4, FOURC, CORE, CORED, COREA)
non hanno segnale a ec20: RMSE di validazione tra 22.23 e 26.39
contro i 12.24 di CORNER, e pesi attivi tra 7 e 349. Le run a
batch pieni lo confermano: RMSE invariate a due decimali e
ottimizzazione che si ferma anticipatamente
(\texttt{LOW\_PROGRESS}). BARBEL fa eccezione: 15.43 a batch 3/6
e 15.24 a batch pieni (37\,317 e 42\,818 pesi attivi) — le
diagonali hanno segnale, ma restano 3.2 dietro CORNER.

Nel modello completo i sei entrano così: il modello 0147 aggiunge
i cinque al 0134 e mette BARBEL al posto di DIAG8, che lo contiene
completamente — quindici pattern, RMSE di validazione 6.49. Per
decomporre i contributi, i pattern sono aggiunti uno alla volta al
0134 (0148--0152) e BARBEL è provato da solo in sostituzione di
DIAG8 (0153). La tabella~\ref{tab:decomposizione} riassume le sei run.

\begin{table}[ht]
\centering
\small
\begin{tabular}{llrrrr}
\toprule
modello & contenuto & RMSE train & RMSE val & $W$ & iterazioni \\
\midrule
0134 & base & 6.30 & 6.57 & 670\,519 & 944 \\
0148 & + RCT2X4 & 6.30 & 6.57 & 670\,869 & 955 \\
0149 & + FOURC & 6.30 & 6.57 & 670\,540 & 1009 \\
0150 & + CORE & 6.30 & 6.57 & 670\,525 & 1007 \\
0151 & + CORED & 6.30 & 6.57 & 670\,750 & 973 \\
0152 & + COREA & 6.30 & 6.57 & 670\,851 & 904 \\
0153 & BARBEL al posto di DIAG8 & 6.22 & 6.50 & 711\,853 & 950 \\
0147 & i cinque + BARBEL insieme & 6.21 & 6.49 & 712\,793 & 1136 \\
\bottomrule
\end{tabular}
\caption{La decomposizione del modello 0147: i cinque pattern
semplici aggiunti uno alla volta al 0134, BARBEL in sostituzione di
DIAG8, e i due estremi.}
\label{tab:decomposizione}
\end{table}

I cinque pattern semplici sono inerti nel modello completo:
aggiunti uno alla volta non muovono la RMSE di validazione (6.57
in tutti e cinque) e portano al 0147 in totale 940 pesi attivi.
La sostituzione BARBEL/DIAG8 da sola vale 0.07 (6.50): è l'intero
miglioramento del 0147 (6.49), al quale i cinque pattern
contribuiscono il residuo 0.01, a livello di risoluzione. Il
modello 0153 — il 0134 con BARBEL al posto di DIAG8 — diventa il
nuovo riferimento della serie: RMSE di training 6.22, di
validazione 6.50, 711\,853 pesi attivi, 950 iterazioni.
